\begin{exercises}
  \recommended \item 
     求下列各矩阵的行列式。
     \begin{exparts*}
       \partsitem \(
            \begin{mat}[r]
                3    &1   \\
               -1    &1
            \end{mat}    \)
       \partsitem \(
            \begin{mat}[r]
                2    &0   &1  \\
                3    &1   &1 \\
               -1    &0   &1
            \end{mat}    \)
       \partsitem \(
            \begin{mat}[r]
                4    &0   &1  \\
                0    &0   &1 \\
                1    &3   &-1
            \end{mat}    \)
     \end{exparts*}
     \begin{answer}
       \begin{exparts*}
         \partsitem \( 4 \)
         \partsitem \( 3 \)
         \partsitem \( -12 \)
       \end{exparts*}  
     \end{answer}
  \item 
     求下列各矩阵的行列式。
     \begin{exparts*}
        \partsitem \( \begin{mat}[r]
                    2  &0  \\
                   -1  &3
                 \end{mat} \)
        \partsitem \( \begin{mat}[r]
                    2  &1  &1  \\
                    0  &5  &-2 \\
                    1  &-3 &4
                 \end{mat} \)
        \partsitem \( \begin{mat}[r]
                    2  &3  &4  \\
                    5  &6  &7  \\
                    8  &9  &1
                 \end{mat} \)
     \end{exparts*}
    \begin{answer}
      \begin{exparts*}
        \partsitem \( 6 \)
        \partsitem \( 21 \)
        \partsitem \( 27 \)
      \end{exparts*}  
    \end{answer}
  \recommended \item  
    验证上三角的
    $\nbyn{3}$ 矩阵的行列式是对角元之积。
    \begin{equation*}
       \det(
       \begin{mat}
           a    &b   &c    \\
           0    &e   &f    \\
           0    &0   &i
       \end{mat}
       )
       =aei
    \end{equation*}
    下三角矩阵也是如此吗？
    \begin{answer}
      对于第一个，应用本节中的公式，注意任何
      含有 \( d \)、\( g \) 或 \( h \) 的项都是零，再化简即可。
      下三角矩阵也是如此。  
    \end{answer}
  \recommended \item 
     用行列式判断下列各矩阵是奇异的还是
     非奇异的。
     \begin{exparts*}
       \partsitem \(
            \begin{mat}[r]
                2    &1   \\
                3    &1
            \end{mat}    \)
       \partsitem \(
            \begin{mat}[r]
                0    &1   \\
                1    &-1
            \end{mat}    \)
       \partsitem \(
            \begin{mat}[r]
                4    &2   \\
                2    &1
            \end{mat}    \)
     \end{exparts*}
     \begin{answer}
       \begin{exparts}
         \partsitem 非奇异的，行列式为 \( -1 \)。
         \partsitem 非奇异的，行列式为 \( -1 \)。
         \partsitem 奇异的，行列式为 \( 0 \)。
       \end{exparts}  
     \end{answer}
  \item 
     奇异的还是非奇异的？
     用行列式来判断。
     \begin{exparts*}
       \partsitem \(
            \begin{mat}[r]
                2    &1   &1  \\
                3    &2   &2 \\
                0    &1   &4
            \end{mat}    \)
       \partsitem \(
            \begin{mat}[r]
                1    &0   &1  \\
                2    &1   &1 \\
                4    &1   &3
            \end{mat}    \)
       \partsitem \(
            \begin{mat}[r]
                2    &1   &0  \\
                3    &-2  &0 \\
                1    &0   &0
            \end{mat}    \)
     \end{exparts*}
     \begin{answer}
      \begin{exparts}
         \partsitem 非奇异的，行列式为 \( 3 \)。
         \partsitem 奇异的，行列式为 \( 0 \)。
         \partsitem 奇异的，行列式为 \( 0 \)。
       \end{exparts}  
     \end{answer}
  \recommended \item
    下面每一对矩阵仅相差一次行变换。
    请利用这次行变换
    来比较 \( \det(A) \) 与 \( \det(B) \)。
     \begin{exparts}
        \partsitem \( A=\begin{mat}[r]
                      1  &2  \\
                      2  &3
                   \end{mat} \), 
               \(  B=\begin{mat}[r]
                      1  &2  \\
                      0  &-1
                   \end{mat}  \)
        \partsitem \( A=\begin{mat}[r]
                      3  &1  &0  \\
                      0  &0  &1  \\
                      0  &1  &2
                   \end{mat} \),
              \(   B=\begin{mat}[r]
                      3  &1  &0  \\
                      0  &1  &2  \\
                      0  &0  &1
                   \end{mat}  \)
        \partsitem \( A=\begin{mat}[r]
                      1  &-1 &3  \\
                      2  &2  &-6 \\
                      1  &0  &4
                   \end{mat} \),
              \(   B=\begin{mat}[r]
                      1  &-1 &3  \\
                      1  &1  &-3 \\
                      1  &0  &4
                   \end{mat}  \)
     \end{exparts}
     \begin{answer}
       \begin{exparts}
         \partsitem \( \det(B)=\det(A) \)，通过 \( -2\rho_1+\rho_2 \)
         \partsitem \( \det(B)=-\det(A) \)，通过 
             \( \rho_2\leftrightarrow\rho_3 \)
         \partsitem \( \det(B)=(1/2)\cdot \det(A) \)，通过 \( (1/2)\rho_2 \)
       \end{exparts}   
     \end{answer}
  \recommended \item 按如下方案求这个 $\nbyn{4}$ 矩阵的行列式：~先做高斯消元，
    并留意行列式在倍加时保持不变，在
    对换两行时变号，在某行倍乘时相应倍乘，
    且阶梯形矩阵的行列式等于
    其对角元之积。
    \begin{equation*}
      \begin{mat}
        1 &2 &0  &2 \\
        2 &4 &1  &0 \\
        0 &0 &-1 &3 \\
        3 &-1 &1 &4
      \end{mat}
    \end{equation*}
    \begin{answer}
      高斯消元法正是这样做的。
      % [ 0 -7  1 -2]
      %  swap row 2 with row 4
      % [ 1  2  0  2]
      % [ 0 -7  1 -2]
      % [ 0  0 -1  3]
      % [ 0  0  1 -4]
      %  take 1 times row 3 plus row 4
      % [ 1  2  0  2]
      % [ 0 -7  1 -2]
      % [ 0  0 -1  3]
      % [ 0  0  0 -1]
      \begin{equation*}
        \begin{mat}
          1 &2 &0  &2 \\
          2 &4 &1  &0 \\
          0 &0 &-1 &3 \\
          3 &-1 &1 &4
        \end{mat}
        \grstep[-3\rho_1+\rho_4]{-2\rho_1+\rho_2}
        \grstep{\rho_2\leftrightarrow\rho_4}
        \grstep{\rho_3+\rho_4}      
        \begin{mat}
          1 &2 &0  &2 \\
          0 &-7 &1  &-2 \\
          0 &0 &-1 &3 \\
          0 &0 &0 &-1
        \end{mat}
      \end{equation*}
      阶梯形矩阵的对角元之积为 
       $1\cdot(-7)\cdot(-1)\cdot(-1)=-7$。
      在高斯消元的过程中没有行被倍乘，但有
      一次行对换，所以求行列式时须变号，
      得 $+7$。
    \end{answer}
  \item 
     证明此式。
      \begin{equation*}
         \det(
         \begin{mat}
             1    &1   &1    \\
             a    &b   &c    \\
             a^2  &b^2 &c^2
         \end{mat}
         )
         =(b-a)(c-a)(c-b)
      \end{equation*}
     \begin{answer}
       利用 $\nbyn{3}$ 矩阵的行列式公式，
       我们将左边展开
       \begin{equation*}
         1\cdot b\cdot c^2+1\cdot c\cdot a^2+1\cdot a\cdot b^2
          -b^2\cdot c\cdot 1 -c^2\cdot a\cdot 1-a^2\cdot b\cdot 1
       \end{equation*}
       而将右边展开则得
       \begin{equation*}
         (bc-ba-ac+a^2)\cdot(c-b)
         =c^2b-b^2c-bac+b^2a-ac^2+acb+a^2c-a^2b
       \end{equation*}
       于是只需验证两者相等即可。
       （\textit{注}。
       这是
       \definend{Vandermonde 行列式}\index{Vandermonde!determinant}%
       \index{determinant!Vandermonde}的 \( \nbyn{3} \) 情形，
       它出现在许多应用中）。
     \end{answer}
   \recommended \item 
      哪些实数 \( x \) 使这个矩阵是奇异的？
      \begin{equation*}
         \begin{mat}
            12-x  &4  \\
            8    &8-x
         \end{mat}
      \end{equation*}
      \begin{answer}
         这个方程
         \begin{equation*}
           0=
           \det(
             \begin{mat}
                12-x  &4  \\
                8    &8-x
             \end{mat}
           )
           =64-20x+x^2
           =(x-16)(x-4) 
       \end{equation*}
       有根 \( x=16 \) 和 \( x=4 \)。  
     \end{answer}
  \item \label{exer:ThreeByThreeDetForm} 
    做高斯消元，以验证本节引言中
    所述的 $\nbyn{3}$ 矩阵公式。
    \begin{center}
      \( \begin{mat}
               a  &b  &c  \\
               d  &e  &f  \\
               g  &h  &i
         \end{mat} \)
      是非奇异的当且仅当
      \( aei+bfg+cdh-hfa-idb-gec \neq 0 \)
    \end{center}
    \begin{answer}
      我们先把矩阵化成阶梯形。
      首先，假设 \( a\neq 0 \) 且 \( ae-bd\neq 0 \)。
      \begin{multline*}
        \grstep{(1/a)\rho_1}
        \begin{mat}
           1   &b/a   &c/a   \\
           d   &e     &f     \\
           g   &h     &i
         \end{mat}                                           
        \grstep[-g\rho_1+\rho_3]{-d\rho_1+\rho_2}
        \begin{mat}
           1   &b/a           &c/a           \\
           0   &(ae-bd)/a     &(af-cd)/a     \\
           0   &(ah-bg)/a     &(ai-cg)/a
         \end{mat}                                            \\
        \grstep{(a/(ae-bd))\rho_2}
        \begin{mat}
           1   &b/a           &c/a             \\
           0   &1             &(af-cd)/(ae-bd) \\
           0   &(ah-bg)/a     &(ai-cg)/a
         \end{mat}
      \end{multline*}
      这一步完成了计算。
      \begin{equation*}
        \grstep{((ah-bg)/a)\rho_2+\rho_3}
        \begin{mat}
           1   &b/a    &c/a             \\
           0   &1      &(af-cd)/(ae-bd)      \\
           0   &0      &(aei+bgf+cdh-hfa-idb-gec)/(ae-bd)
         \end{mat}
      \end{equation*}
      现在假设 $a\neq 0$ 且 \( ae-bd\neq 0 \)，
      原矩阵是非奇异的
      当且仅当上面第 \( 3,3 \) 处的元素非零。
      也就是说，在这些假设下，原矩阵
      是非奇异的当且仅当 $aei+bgf+cdh-hfa-idb-gec\neq 0$，
      即为所证。

      最后我们逐一看一下，若为方便起见而在前一段中所作的假设不成立，
      会发生什么。
      首先，若 \( a\neq 0 \) 但 \( ae-bd=0 \)，那么我们可以对换
      \begin{equation*}
        \begin{mat}
           1   &b/a           &c/a           \\
           0   &0             &(af-cd)/a     \\
           0   &(ah-bg)/a     &(ai-cg)/a
         \end{mat}                                  
        \grstep{\rho_2\leftrightarrow\rho_3}
        \begin{mat}
           1   &b/a           &c/a           \\
           0   &(ah-bg)/a     &(ai-cg)/a     \\
           0   &0             &(af-cd)/a
         \end{mat}
      \end{equation*}
      从而得出结论：矩阵是非奇异的当且仅当
      \( ah-bg=0 \) 或 \( af-cd=0 \)。
      条件‘\( ah-bg=0 \) 或 \( af-cd=0 \)’等价于
      条件‘\( (ah-bg)(af-cd)=0 \)’。
      乘开并利用本情形的假设 $ae-bd=0$，以 $ae$ 代 $bd$，便得
      \begin{multline*}
         0=ahaf-ahcd-bgaf+bgcd
          =ahaf-ahcd-bgaf+aegc            \\
          =a(haf-hcd-bgf+egc)
      \end{multline*}
      由于 \( a\neq 0 \)，故矩阵
      是非奇异的当且仅当 \( haf-hcd-bgf+egc=0 \)。
      因此，在 \( a\neq 0 \) 且 \( ae-bd=0 \) 的情形下，
      当
      \( haf-hcd-bgf+egc-i(ae-bd)=0 \) 时矩阵
      是非奇异的。

      余下的情形是常规的。
      对于 \( a=0 \) 但 \( d\neq 0 \) 的情形，以及 \( a=0 \)、\( d=0 \)
      但 \( g\neq 0 \) 的情形，都先对换两行，然后
      照上面那样继续即可。
      \( a=0 \)、\( d=0 \) 且 \( g=0 \) 的情形则很简单\Dash 该矩阵是
      奇异的，因为它的列构成一个线性相关的组，而且
      行列式算出来为零。  
    \end{answer}
  \item 
     证明 \( \Re^2 \) 中过 \( (x_1,y_1) \)
     与 \( (x_2,y_2) \) 的直线方程由这个行列式给出。
     \begin{equation*}
        \det(
        \begin{mat}
           x   &y   &1  \\
           x_1 &y_1 &1  \\
           x_2 &y_2 &1
        \end{mat})=0 \qquad x_1\neq x_2
     \end{equation*}
     \begin{answer}
       算出行列式并做一点代数运算，便得到
       \begin{align*}
          0
          &=y_1x+x_2y+x_1y_2-y_2x-x_1y-x_2y_1     \\
          (x_2-x_1)\cdot y
          &=(y_2-y_1)\cdot x+x_2y_1-x_1y_2              \\
          y
          &=\frac{y_2-y_1}{x_2-x_1}\cdot x+\frac{x_2y_1-x_1y_2}{x_2-x_1}
       \end{align*}
       注意这是一个直线的方程（特别地，
       其中含有熟悉的斜率表达式），
       并注意 \( (x_1,y_1) \) 与 \( (x_2,y_2) \) 都满足它。 
     \end{answer}
  \item
    许多人学过
    计算这个 \( \nbyn{3} \)
    矩阵行列式的口诀：~把前两列抄到矩阵右侧，
    然后取沿正向对角线的乘积并相加，
    再取沿反向对角线的乘积并相减。
    也就是说，先写成
    \begin{equation*}
        \begin{pmat}{ccc|cc}
          h_{1,1} &h_{1,2} &h_{1,3} &h_{1,1} &h_{1,2} \\
          h_{2,1} &h_{2,2} &h_{2,3} &h_{2,1} &h_{2,2} \\
          h_{3,1} &h_{3,2} &h_{3,3} &h_{3,1} &h_{3,2}
        \end{pmat}
     \end{equation*}
     然后计算这个式子。
     \begin{equation*}
      \begin{array}{l}
      h_{1,1}h_{2,2}h_{3,3}+h_{1,2}h_{2,3}h_{3,1}+h_{1,3}h_{2,1}h_{3,2} \\
      \>-h_{3,1}h_{2,2}h_{1,3}-h_{3,2}h_{2,3}h_{1,1}
        -h_{3,3}h_{2,1}h_{1,2}
      \end{array}
    \end{equation*}
    \begin{exparts}
      \partsitem 验证这与本节引言中
        给出的公式一致。
      \partsitem 它能否推广到其他阶数的行列式？
    \end{exparts}
    \begin{answer}
      \begin{exparts}
        \partsitem 与本节引言中给出的公式比较，很容易。
        \partsitem 它对 \( \nbyn{2} \) 矩阵成立
          \begin{align*}
            \begin{pmat}{cc|c}
              h_{1,1} &h_{1,2} &h_{1,1} \\
              h_{2,1} &h_{2,2} &h_{2,1}
            \end{pmat}
            &=\begin{array}[t]{@{}l@{}}
              h_{1,1}h_{2,2}+h_{1,2}h_{2,1}  \\
              \>-h_{2,1}h_{1,2}-h_{2,2}h_{1,1}
            \end{array}                                     \\
            &=h_{1,1}h_{2,2}-h_{1,2}h_{2,1}
          \end{align*}
          但对 \( \nbyn{4} \) 矩阵不成立。
          例如，这个矩阵是
          奇异的，因为第二行与第三行相等
          \begin{equation*}
            \begin{mat}[r]  
              1  &0  &0  &1   \\
              0  &1  &1  &0   \\
              0  &1  &1  &0   \\
             -1  &0  &0  &1  
            \end{mat}
          \end{equation*}
          但按口诀的方案算并不得到零。
          \begin{equation*}
            \begin{pmat}{rrrr|rrr}
                       1  &0  &0  &1  &1  &0  &0  \\
                       0  &1  &1  &0  &0  &1  &1  \\
                       0  &1  &1  &0  &0  &1  &1  \\
                      -1  &0  &0  &1  &-1 &0  &0
            \end{pmat}
            =\begin{array}[t]{@{}l@{}}
               1+0+0+0 \\
               \>-(-1)-0-0-0
              \end{array}
           \end{equation*}
      \end{exparts}  
     \end{answer}
  \item 
    向量
    \begin{equation*}
      \vec{x}=\colvec{x_1 \\ x_2 \\ x_3}
      \qquad
      \vec{y}=\colvec{y_1 \\ y_2 \\ y_3}
    \end{equation*}
    的
    \definend{叉积}\index{cross product}\index{vector!cross product} 
    是按这个行列式算出的向量。
    \begin{equation*}
      \vec{x}\times\vec{y}=
      \det(\begin{mat}
        \vec{e}_1  &\vec{e}_2  &\vec{e}_3  \\
        x_1        &x_2        &x_3        \\
        y_1        &y_2        &y_3
      \end{mat})
    \end{equation*}
    注意第一行的元素是向量，即 $\Re^3$ 的
    标准基中的向量。
    证明两个向量的叉积与每个向量都垂直。
    \begin{answer}
      该行列式为 
      $
        (x_2y_3-x_3y_2)\vec{e}_1
          +(x_3y_1-x_1y_3)\vec{e}_2
          +(x_1y_2-x_2y_1)\vec{e}_3
      $。
      要检验垂直性，我们检验它与第一个向量的点积
      为零
      \begin{equation*}
        \colvec{x_1 \\ x_2 \\ x_3}
        \dotprod
        \colvec{x_2y_3-x_3y_2 \\ x_3y_1-x_1y_3 \\ x_1y_2-x_2y_1}
        =x_1x_2y_3-x_1x_3y_2+x_2x_3y_1-x_1x_2y_3+x_1x_3y_2-x_2x_3y_1=0
      \end{equation*}
      而它与第二个向量的点积也为零。
      \begin{equation*}
        \colvec{y_1 \\ y_2 \\ y_3}
        \dotprod
        \colvec{x_2y_3-x_3y_2 \\ x_3y_1-x_1y_3 \\ x_1y_2-x_2y_1}
        =x_2y_1y_3-x_3y_1y_2+x_3y_1y_2-x_1y_2y_3+x_1y_2y_3-x_2y_1y_3=0
      \end{equation*}  
    \end{answer}
  \item 
    证明下列每个论断对 $\nbyn{2}$ 矩阵成立。  
    \begin{exparts}
      \partsitem 乘积的行列式
        等于行列式的乘积
        $\det(ST)=\det(S)\cdot\det(T)$。
      \partsitem 若 \( T \) 是可逆的，则
        逆矩阵的行列式等于行列式的倒数
        \( \det(T^{-1})=(\,\det(T)\,)^{-1} \)。
    \end{exparts}
    矩阵 $T$ 与 $T^\prime$ 称为
    \definend{相似的}\index{similar}，若存在
    非奇异矩阵 $P$ 使得 $T^\prime=PTP^{-1}$。
    （我们将在第五章考察这一关系。）
    证明相似的 \( \nbyn{2} \) 矩阵有相同的
    行列式。
    \begin{answer}
       \begin{exparts}
        \partsitem 直接代入计算：
          乘积的行列式为
          \begin{align*}
             \det(\begin{mat}
                       a  &b  \\
                       c  &d
                    \end{mat}
                    \begin{mat}
                       w  &x  \\
                       y  &z
                    \end{mat}  )
             &=
             \det(\begin{mat}
                 aw+by  &ax+bz  \\
                 cw+dy  &cx+dz
              \end{mat} )                 \\
             &=
             \begin{array}[t]{@{}l@{}} 
                 acwx+adwz+bcxy+bdyz  \\
                 \> -acwx-bcwz-adxy-bdyz
             \end{array}
          \end{align*}
          而行列式的乘积为
          \begin{equation*}
             \det(\begin{mat}
                a  &b  \\
                c  &d
             \end{mat})
             \cdot\det(\begin{mat}
                w  &x  \\
                y  &z
             \end{mat})
             =
             (ad-bc)\cdot (wz-xy)
          \end{equation*}
          验证两者相等很容易。
        \partsitem 利用前一小题。
       \end{exparts}  
       \noindent 由上面两条立即得到：相似矩阵有相同的行列式
       $
          \det(PTP^{-1})=\det(P)\cdot\det(T)\cdot\det(P^{-1})
       $。
     \end{answer}
  \recommended \item
    证明平面内这个区域的面积
    \begin{center}
      \includegraphics{det/mp/ch4.30}
    \end{center}
    等于这个行列式的值。
    \begin{equation*}
       \det(
       \begin{mat}
          x_1  &x_2  \\
          y_1  &y_2
       \end{mat})
    \end{equation*}
    并与这个比较。
    \begin{equation*}
       \det(
       \begin{mat}
          x_2  &x_1  \\
          y_2  &y_1
       \end{mat})
    \end{equation*}
    \begin{answer}
      一种方法是数出这些面积
      \begin{center}
        \includegraphics{det/mp/ch4.31}
      \end{center}
      即取整个矩形的面积，再减去
      左上矩形 $A$、上中三角形 $B$、
      右上三角形 $D$、左下三角形 $C$、
      下中三角形 $E$ 与右下矩形 $F$ 的面积
      \( (x_1+x_2)(y_1+y_2)-x_2y_1-(1/2)x_1y_1-(1/2)x_2y_2
              -(1/2)x_2y_2-(1/2)x_1y_1-x_2y_1 \)。
      化简即得行列式公式。

      这个行列式是上面那个的相反数；该公式
      区分了第二列是否在第一列的
      逆时针方向。  
     \end{answer}
  \item 
    证明对 \( \nbyn{2} \) 矩阵，矩阵的行列式
    等于其转置的行列式。
    这对 \( \nbyn{3} \) 矩阵也成立吗？
    \begin{answer}
      对 \( \nbyn{2} \) 矩阵，利用
      引言中引用的公式，计算很容易。
      它对 \( \nbyn{3} \) 矩阵也成立；其
      计算是常规的。  
    \end{answer}
  \item 
    行列式函数是线性的吗 \Dash  即
    \( \det(x\cdot T+y\cdot S)=x\cdot \det(T)+y\cdot \det(S) \) 成立吗？
    \begin{answer}
      不成立。
      我们用 $\nbyn{2}$ 行列式来说明。
      回想常数可以每次提出一行。
      \begin{equation*}
         \det(
         \begin{mat}[r]
            2  &4  \\
            2  &6  \\
         \end{mat})
         =
         2\cdot\det(\begin{mat}[r]
            1  &2  \\
            2  &6  \\
         \end{mat})
         =
         2\cdot 2\cdot \det(\begin{mat}[r]
            1  &2  \\
            1  &3  \\
         \end{mat})
      \end{equation*}
      这与线性性矛盾（这里我们不需要 \( S \)，也就是说，我们可以
      取 $S$ 为零矩阵）。  
    \end{answer}
  \item 
    证明：若 \( A \) 是 \( \nbyn{3} \) 的，则
    \( \det(c\cdot A)=c^3\cdot \det(A) \) 对任意标量 \( c \) 成立。
    \begin{answer}
       把 \( c \) 一次一行地提出来。  
    \end{answer}
  \item 
    哪些实数 \( \theta \) 使
    \begin{equation*}
       \begin{mat}
          \cos\theta  &-\sin\theta  \\
          \sin\theta  &\cos\theta
       \end{mat}
    \end{equation*}
    奇异？
    从几何上解释。
    \begin{answer}
      不存在使该矩阵奇异的实数 \( \theta \)，
      因为矩阵的行列式
      \( \cos^2\theta+\sin^2\theta \) 永不为 $0$，它对所有 $\theta$
      都等于 $1$。
      从几何上看，关于标准基，
      这个矩阵表示
      平面旋转角度 \( \theta \) 的转动。
      每个这样的映射都是单射 \Dash  例如，它是可逆的。  
    \end{answer}
  \puzzle \item  
    \cite{Monthly55p257}
    若一个三阶行列式的元素为
    \( 1 \), \( 2 \), \ldots, \( 9 \)，它可能取到的最大值
    是多少？
    \begin{answer}
      \answerasgiven
      设 \( P \) 为行列式中三个正项之和，
      \( -N \) 为三个负项之和。
      \( P \) 的最大值为
      \begin{equation*}
        9\cdot 8\cdot 7 +6\cdot 5\cdot 4 +3\cdot 2\cdot 1=630.
      \end{equation*}
      与 \( P \) 相容的 \( N \) 的最小值为
      \begin{equation*}
        9\cdot 6\cdot 1 +8\cdot 5\cdot 2 +7\cdot 4\cdot 3=218.
      \end{equation*}
      \( P \) 的任何改变都会使该和的降低量超过
      \( 4 \)。
      因此 \( 412 \) 是该行列式的最大值，而该行列式的一种形式是
      \begin{equation*}
         \begin{vmat}[r]
            9  &4  &2  \\
            3  &8  &6  \\
            5  &1  &7
         \end{vmat}.
      \end{equation*}  
    \end{answer}
\end{exercises}
